\section{Induction}

\begin{frame}
\centering
\Huge\textbf{Induction}
\end{frame}

% --------------------------------------------------
% Question 1
% --------------------------------------------------

\begin{frame}{Question 1}
A program processes $n$ records. The number of elementary processing steps is modeled by

$$
T(n)=1+2+\cdots+n.
$$

A developer claims that

$$
T(n)=\frac{n(n+1)}{2}
$$

for every positive integer $n$.

\textbf{Question:} Prove the claim using mathematical induction.
\end{frame}

\begin{frame}{Answer 1}
Define the statement

$$
P(n):\quad
1+2+\cdots+n=\frac{n(n+1)}{2}.
$$

\textbf{Base case:} $n=1$

$$
1=\frac{1(1+1)}{2}=1.
$$

Therefore, $P(1)$ is true.
\end{frame}

\begin{frame}{Answer 1}
\textbf{Inductive hypothesis:}

Assume $P(k)$ is true:

$$
1+2+\cdots+k
=
\frac{k(k+1)}{2}.
$$

We must prove $P(k+1)$:

$$
1+2+\cdots+k+(k+1)
=
\frac{(k+1)(k+2)}{2}.
$$

\end{frame}

\begin{frame}{Answer 1}
Using the inductive hypothesis,

$$
1+2+\cdots+k+(k+1)
$$
$$
=
\frac{k(k+1)}{2}+(k+1)
$$
$$
=
\frac{k(k+1)+2(k+1)}{2}
$$
$$
=
\frac{(k+1)(k+2)}{2}.
$$
Therefore, $P(k+1)$ is true.
Hence,
$$
\boxed{
1+2+\cdots+n=\frac{n(n+1)}{2}
}
$$

for every positive integer $n$.
\end{frame}

% --------------------------------------------------
% Question 2
% --------------------------------------------------

\begin{frame}{Question 2}
A data structure stores a sequence of nodes. Suppose the number of
operations required to process the first $n$ nodes is

$$
T(n)=1+3+5+\cdots+(2n-1).
$$

\textbf{Question:} Prove using mathematical induction that

$$
T(n)=n^2
$$

for every positive integer $n$.
\end{frame}

\begin{frame}{Answer 2}
Define

$$
P(n):\quad
1+3+5+\cdots+(2n-1)=n^2.
$$

\textbf{Base case:} $n=1$

$$
1=1^2.
$$

Therefore, $P(1)$ is true.
\end{frame}

\begin{frame}{Answer 2}
\textbf{Inductive hypothesis:}

Assume

$$
1+3+5+\cdots+(2k-1)=k^2.
$$

We must prove

$$
1+3+5+\cdots+(2k-1)+(2k+1)
=
(k+1)^2.
$$

Using the hypothesis,

$$
k^2+(2k+1)
=
k^2+2k+1
=
(k+1)^2.
$$

\end{frame}

\begin{frame}{Answer 2}
Therefore,

$$
1+3+5+\cdots+(2k+1)
=
(k+1)^2.
$$

Thus, $P(k+1)$ is true.

Hence, by mathematical induction,

$$
\boxed{
1+3+5+\cdots+(2n-1)=n^2
}
$$

for every positive integer $n$.
\end{frame}

% --------------------------------------------------
% Question 3
% --------------------------------------------------

\begin{frame}{Question 3}
A data structure contains a binary tree with $n$ levels. Suppose the maximum
number of nodes at level $i$ is

$$
2^{i-1}.
$$

\textbf{Question:} Prove that the maximum total number of nodes in a binary
tree with $n$ levels is

$$
2^n-1.
$$

\end{frame}

\begin{frame}{Answer 3}
Let

$$
P(n):\quad
1+2+2^2+\cdots+2^{n-1}=2^n-1.
$$

\textbf{Base case:} $n=1$

$$
1=2^1-1=1.
$$

Therefore, $P(1)$ is true.
\end{frame}

\begin{frame}{Answer 3}
Assume

$$
1+2+\cdots+2^{k-1}=2^k-1.
$$

For $k+1$ levels,

$$
1+2+\cdots+2^{k-1}+2^k.
$$

Using the hypothesis,

$$
=(2^k-1)+2^k
$$

$$
=2^{k+1}-1.
$$

Therefore, $P(k+1)$ is true.
\end{frame}

\begin{frame}{Answer 3}
Hence, by mathematical induction,

$$
\boxed{
1+2+2^2+\cdots+2^{n-1}=2^n-1
}
$$

Therefore, the maximum number of nodes in a binary tree with $n$ levels is

$$
\boxed{2^n-1}.
$$

\end{frame}

% --------------------------------------------------
% Question 4
% --------------------------------------------------

\begin{frame}{Question 4}
A university maintains three sets of students:

$$
A=\text{students enrolled in Artificial Intelligence},
$$

$$
B=\text{students enrolled in Data Science},
$$

and

$$
C=\text{students enrolled in Cybersecurity}.
$$

To avoid counting the same student multiple times, the university uses the
inclusion--exclusion principle.

\textbf{Question:} Prove the three-set inclusion--exclusion formula using
mathematical induction as a step toward the general principle.
\end{frame}

\begin{frame}{Answer 4}
For three sets, the inclusion--exclusion formula is

$$
|A\cup B\cup C|
=
|A|+|B|+|C|
-|A\cap B|-|A\cap C|-|B\cap C|
+|A\cap B\cap C|.
$$

For the general principle, define

$$
P(n):
\left|
\bigcup_{i=1}^{n}A_i
\right|
=
\sum_i|A_i|
-\sum_{i<j}|A_i\cap A_j|
+\sum_{i<j<k}|A_i\cap A_j\cap A_k|
-\cdots.
$$

\end{frame}

\begin{frame}{Answer 4}
The formula follows by counting each element according to the number of
sets containing it.

Suppose an element belongs to exactly $r$ of the sets.

It is counted:
$$
\binom{r}{1}
$$

times in the single-set terms,
$$
\binom{r}{2}
$$

times in the pairwise terms, and so on.

Therefore, its total contribution is
$$
\binom{r}{1}
-\binom{r}{2}
+\binom{r}{3}
-\cdots
+(-1)^{r+1}\binom{r}{r}.
$$

\end{frame}

\begin{frame}{Answer 4}
Using the binomial identity

$$
(1-1)^r
=
\binom{r}{0}
-\binom{r}{1}
+\binom{r}{2}
-\cdots
+(-1)^r\binom{r}{r},
$$

we obtain

$$
\binom{r}{1}
-\binom{r}{2}
+\cdots
+(-1)^{r+1}\binom{r}{r}
=1.
$$

Thus, every element belonging to the union is counted exactly once.

Hence,

$$
\boxed{
\left|\bigcup_{i=1}^{n}A_i\right|
=
\sum_{k=1}^{n}
(-1)^{k+1}
\sum_{1\le i_1<\cdots<i_k\le n}
|A_{i_1}\cap\cdots\cap A_{i_k}|
}
$$

\end{frame}

% --------------------------------------------------
% Question 5
% --------------------------------------------------

\begin{frame}{Question 5}
Consider a recursive algorithm defined by

$$
F(1)=1
$$

and

$$
F(n)=F(n-1)+2
\qquad(n>1).
$$

A programmer claims that

$$
F(n)=2n-1.
$$

\textbf{Question:} Prove the claim using mathematical induction.
\end{frame}

\begin{frame}{Answer 5}
Let

$$
P(n):\quad F(n)=2n-1.
$$

\textbf{Base case:}

$$
F(1)=1
$$

and

$$
2(1)-1=1.
$$

Therefore, $P(1)$ is true.
\end{frame}

\begin{frame}{Answer 5}
Assume

$$
F(k)=2k-1.
$$

Using the recursive definition,

$$
F(k+1)=F(k)+2.
$$

By the inductive hypothesis,

$$
F(k+1)
=
(2k-1)+2
=
2k+1.
$$

Therefore,

$$
F(k+1)=2(k+1)-1.
$$

\end{frame}

\begin{frame}{Answer 5}
Thus, $P(k+1)$ is true.

Since the base case is true and the inductive step is valid,

$$
\boxed{F(n)=2n-1}
$$

for every positive integer $n$.
\end{frame}
